\documentclass[11pt]{article}
\usepackage[hyperref]{acl}
\usepackage{times}
\usepackage{latexsym}
\usepackage{graphicx}
\usepackage{booktabs}
\usepackage{amsmath}
\usepackage{amssymb}
\usepackage{multirow}
\usepackage{xcolor}
\usepackage{enumitem}
\usepackage{url}

\title{Differential Depth Response in Recursive LLM Generation:\\Classifying Generation Depth Beyond Binary Detection}

\author{Anonymous}

\begin{document}
\maketitle

%=====================================================================
\begin{abstract}
Web-crawled training corpora increasingly contain LLM-generated text of unknown provenance depth: text may be human-authored (gen-0), directly generated (gen-1), or the product of recursive generation chains (gen-2+). Current detection methods offer only binary real-vs-synthetic classification, conflating qualitatively different material. We formulate \textbf{generation depth classification} as a new task and characterize \textbf{differential depth response} (DDR): central-tendency features (entropy mean, TTR, hapax ratio) drift monotonically with depth while shape features (entropy max, tail decay) remain near-stationary. On 5 base-model generators from 3 families (0.6B--8B), a 7-feature distributional ensemble achieves leave-one-generator-out (LOGO) binary F1 = 0.680 and macro F1 = 0.762 under 5-fold cross-validation. Cross-model chains yield binary F1 = 0.764. We report a negative finding: depth-aware data curation is statistically equivalent to random subsampling for downstream CLM quality (TOST $p = 1.1 \times 10^{-4}$ at $\pm 2\%$ margin), positioning depth fingerprinting as a diagnostic tool for provenance characterization rather than a curative tool for data selection. The depth signal degrades at larger model scales (8B: $-10.6$ pp vs.\ average), under instruction tuning, and with human editing ($-5.3$ pp for light rephrasing), defining scope boundaries for the approach.
\end{abstract}

%=====================================================================
\section{Introduction}
\label{sec:intro}

The proliferation of large language model (LLM) outputs on the web has created a data provenance crisis. Estimates suggest that 7--17\% of recent scientific papers contain substantial AI-generated content, and web-crawled corpora used for pretraining increasingly include text that was itself generated by earlier models---creating recursive generation chains of unknown depth.

Current AI-text detection methods frame the problem as binary classification: human-written vs.\ machine-generated. Detectors such as FastDetectGPT \cite{bao2024fastdetectgpt}, Binoculars \cite{hans2024binoculars}, and LD-Score provide real-vs-synthetic decisions but cannot distinguish first-generation output (gen-1) from deeper recursive generations (gen-2+). This distinction matters because recursive generation compounds distributional drift: vocabulary diversity decreases, perplexity compresses, and lexical repetition increases with each generation step.

We formulate \textbf{generation depth classification} as a three-class problem: gen-0 (human), gen-1 (first generation), and gen-2+ (recursive generation). To motivate this formulation, we characterize \textbf{differential depth response} (DDR): the empirical observation that distributional features respond differently to increasing generation depth. Specifically, central-tendency and diversity features (entropy mean, TTR, hapax ratio) exhibit monotonic drift with depth ($|\rho| > 0.55$, $R^2 > 0.79$), while distributional shape features (entropy max, tail decay) remain near-stationary ($|\rho| < 0.10$, $R^2 < 0.07$). This differential pattern is consistent with recursive generation inducing a location shift in the NLL distribution while approximately preserving distributional shape.

Our contributions are:
\begin{enumerate}[nosep]
    \item \textbf{Task formulation}: Generation depth classification beyond binary detection, with a 7-feature distributional ensemble achieving LOGO binary F1 = 0.680 across 5 generators from 3 families.
    \item \textbf{DDR characterization}: Empirical demonstration that drifting and stationary features capture complementary depth signal, validated across alternative reference models.
    \item \textbf{Negative finding}: Depth-aware data curation does not improve downstream CLM quality (TOST confirmed), establishing depth fingerprinting as a diagnostic rather than curative tool.
\end{enumerate}

We evaluate across 13 experimental conditions covering feature ablation, cross-generator generalization, cross-family transfer, alternative reference models, prompt-type dependence, paraphrastic chains, cross-model chains, multi-domain seeds, finer-grained ordinal classification, human edit robustness, character-level features, end-to-end neural classifiers, text length confounds, larger-scale generators, and ordinal baseline comparisons.

%=====================================================================
\section{Related Work}
\label{sec:related}

\paragraph{AI-Text Detection.}
Statistical detection methods leverage distributional properties of model-generated text. DetectGPT \cite{mitchell2023detectgpt} uses perturbation-based log-probability curvature; FastDetectGPT \cite{bao2024fastdetectgpt} replaces perturbations with conditional probability estimation. Binoculars \cite{hans2024binoculars} computes cross-perplexity ratios between two models. GLTR \cite{gehrmann2019gltr} visualizes per-token likelihood rankings. DNA-GPT introduces divergent n-gram analysis. These methods provide binary scores but do not estimate generation depth.

\paragraph{Model Collapse and Iterative Generation.}
Shumailov et al.\ (2024) demonstrate model collapse when training on self-generated data across generations. Briesch et al.\ (2023) study large language models as ``broken telephones,'' characterizing information degradation in iterative generation chains. Guo et al.\ (2024) formalize attractor convergence in paraphrastic generation chains, bounding effective diversity to $\sim$3 classes. Our work extends this line by characterizing non-paraphrastic continuation chains, where attractor dynamics differ.

\paragraph{Data Quality and Provenance.}
Penedo et al.\ (2023) develop FineWeb quality scoring for pretraining data. Lee et al.\ (2022) study deduplication for training data quality. Our negative finding on depth-aware curation connects to this literature: knowing generation depth does not provide actionable signal for data selection beyond binary filtering.

\paragraph{Vocabulary Diversity and Stylometry.}
TTR, hapax ratio, and related lexical diversity metrics have long been used in stylometry and authorship attribution. Our surface features draw directly from this tradition. The connection between generation and vocabulary impoverishment is well-established; our contribution is the ordinal depth estimation formulation and the complementarity between surface and reference features.

%=====================================================================
\section{Method}
\label{sec:method}

\subsection{Problem Formulation}

Given a text passage $x$, we classify its \textbf{generation depth} $d \in \{0, 1, 2+\}$, where $d=0$ denotes human-authored text, $d=1$ denotes first-generation LLM output from a human seed, and $d=2+$ denotes text produced by recursively feeding LLM output back as input. The primary evaluation metric is \textbf{leak-free binary F1} (gen-1 vs.\ gen-2+), which excludes gen-0 to avoid seed-sharing leakage between training and test sets.

\subsection{Generation Chain Construction}

We construct recursive generation chains by: (1) sampling a human-authored seed text from a source corpus, (2) generating a continuation using an LLM (gen-1), (3) feeding the generated text back as input to produce gen-2, and repeating up to depth 4. All generations use temperature $t = 0.7$ and top-$p = 0.95$ by default. We use three prompt types: \textit{continuation} (direct text completion), \textit{elaboration} (expand on the given text), and \textit{inspired rewriting} (write a new text inspired by the input).

\subsection{Feature Design}
\label{sec:features}

We extract 7 distributional features from each text, organized into two categories based on their depth response behavior:

\paragraph{Surface Features (2).} Computed directly from the text: \textbf{Type-Token Ratio (TTR)}: ratio of unique tokens to total tokens; \textbf{Hapax Ratio}: proportion of tokens appearing exactly once.

\paragraph{Reference Features (5).} Computed using a fixed reference model (Pythia-410M by default). For each token $x_t$ in the text, we obtain the reference model's predictive distribution $P_{\text{ref}}(\cdot | x_{<t})$ and compute the negative log-likelihood $\ell_t = -\log P_{\text{ref}}(x_t | x_{<t})$. The 5 features are: \textbf{entropy\_mean}: mean of $\ell_t$; \textbf{entropy\_std}: standard deviation of $\ell_t$; \textbf{entropy\_min}: minimum $\ell_t$; \textbf{tail\_decay}: rate of exponential tail decay in the $\ell_t$ distribution; \textbf{entropy\_max}: maximum $\ell_t$.

\paragraph{DDR Categorization.} Features are categorized by their depth-trend behavior:
\begin{itemize}[nosep]
    \item \textbf{Drifting}: monotonic trend with depth ($R^2 > 0.5$)---hapax ratio, TTR, entropy\_mean
    \item \textbf{Stationary}: no systematic trend ($R^2 < 0.1$)---tail\_decay, entropy\_max
    \item \textbf{Ambiguous}: intermediate ($0.1 < R^2 < 0.5$)---entropy\_std, entropy\_min
\end{itemize}

\subsection{Classification}

We use XGBoost as the classifier with 5-fold cross-validation. For cross-generator evaluation, we employ leave-one-generator-out (LOGO) cross-validation: the classifier is trained on all generators except one and tested on the held-out generator. For stricter evaluation, leave-one-family-out (LOFO) holds out all generators from a model family.

%=====================================================================
\section{Experiments}
\label{sec:experiments}

\subsection{Data}

Our primary dataset (MVP) uses Pythia-2.8B with 1000 Wikipedia seeds $\times$ 5 depths, yielding $\sim$4600 records after filtering. The expanded dataset includes 5 generators (Pythia-2.8B, Pythia-1B, OLMo-2-1B, Qwen3-1.7B, Qwen3-0.6B) with 1000 seeds $\times$ 5 generators $\times$ 5 depths $\approx$ 22,800 chains.

\subsection{Feature Ablation and Depth Trends}
\label{sec:ablation}

Table~\ref{tab:ablation} shows the feature ablation results on the MVP dataset. The full 7-feature set achieves macro F1 = 0.773 under 5-fold CV. Surface-only features (TTR + hapax) achieve 0.741, while reference-only features achieve 0.698. The surface-vs-full gap is 3.2 pp (McNemar $p = 0.027$).

\begin{table}[t]
\centering
\small
\caption{Feature ablation: 5-fold CV macro F1 on Pythia-2.8B MVP data.}
\label{tab:ablation}
\begin{tabular}{lccc}
\toprule
\textbf{Feature Set} & \textbf{Macro F1} & \textbf{95\% CI} & \textbf{$\Delta$ vs Full} \\
\midrule
Full (7 features) & 0.773 & [0.758, 0.788] & --- \\
Surface-only (TTR + hapax) & 0.741 & [0.724, 0.758] & $-3.2$ pp \\
Reference-only (5 entropy) & 0.698 & [0.679, 0.717] & $-7.5$ pp \\
\bottomrule
\end{tabular}
\end{table}

Table~\ref{tab:trends} presents per-feature depth trends. Three features exhibit clear monotonic drift with depth: hapax ratio (slope $= -0.063$, $R^2 = 0.88$), TTR (slope $= -0.041$, $R^2 = 0.82$), and entropy\_mean (slope $= -0.089$, $R^2 = 0.79$). In contrast, tail\_decay ($R^2 = 0.04$) and entropy\_max ($R^2 = 0.07$) show no systematic depth trend. Perplexity drops steeply across depths: from 14.2 (depth 0) to 3.8 (depth 2) to 1.9 (depth 4).

\begin{table}[t]
\centering
\small
\caption{Per-feature depth trend (linear regression, depth 0--4).}
\label{tab:trends}
\begin{tabular}{llccc}
\toprule
\textbf{Feature} & \textbf{Category} & \textbf{Slope} & \textbf{$R^2$} & \textbf{$\rho$} \\
\midrule
hapax\_ratio & Drifting & $-0.063$ & 0.88 & $-0.63$ \\
TTR & Drifting & $-0.041$ & 0.82 & $-0.61$ \\
entropy\_mean & Drifting & $-0.089$ & 0.79 & $-0.60$ \\
entropy\_std & Ambiguous & $-0.031$ & 0.41 & $-0.20$ \\
entropy\_min & Stationary & $-0.018$ & 0.22 & $-0.15$ \\
tail\_decay & Stationary & $+0.003$ & 0.04 & $+0.04$ \\
entropy\_max & Stationary & $+0.008$ & 0.07 & $+0.05$ \\
\bottomrule
\end{tabular}
\end{table}

\subsection{Cross-Generator Evaluation}
\label{sec:cross_gen}

Table~\ref{tab:logo} presents leave-one-generator-out (LOGO) results across 5 generators. Average binary F1 is 0.680 with per-generator range of 8.1 pp (OLMo-2-1B: 0.643 to Pythia-1B: 0.724). The within-family inflation effect is evident: Pythia-1B achieves the highest held-out F1 (0.724) because Pythia-2.8B provides same-family signal in the training set.

\begin{table}[t]
\centering
\small
\caption{Leave-one-generator-out binary F1 (gen-1 vs.\ gen-2+).}
\label{tab:logo}
\begin{tabular}{llcc}
\toprule
\textbf{Held-out} & \textbf{Family} & \textbf{Binary F1} & \textbf{95\% CI} \\
\midrule
Pythia-2.8B & Pythia & 0.711 & [0.693, 0.729] \\
Pythia-1B & Pythia & 0.724 & [0.707, 0.741] \\
OLMo-2-1B & OLMo & 0.643 & [0.621, 0.665] \\
Qwen3-1.7B & Qwen & 0.668 & [0.648, 0.688] \\
Qwen3-0.6B & Qwen & 0.655 & [0.634, 0.676] \\
\midrule
\textbf{Average} & --- & \textbf{0.680} & [0.672, 0.688] \\
\bottomrule
\end{tabular}
\end{table}

Leave-one-family-out (LOFO) evaluation further reduces performance (Table~\ref{tab:lofo}). Average LOFO binary F1 = 0.665, with Pythia held-out yielding 0.621 and OLMo held-out yielding 0.726. The within-family inflation from LOGO to LOFO is approximately 3.8 pp for Pythia.

\begin{table}[t]
\centering
\small
\caption{Leave-one-family-out binary F1.}
\label{tab:lofo}
\begin{tabular}{lcc}
\toprule
\textbf{Held-out Family} & \textbf{Binary F1} & \textbf{$\Delta$ vs LOGO avg} \\
\midrule
Pythia (both) & 0.621 & $-5.9$ pp \\
Qwen (both) & 0.648 & $-3.2$ pp \\
OLMo & 0.726 & $+4.6$ pp \\
\midrule
\textbf{Average} & \textbf{0.665} & $-1.5$ pp \\
\bottomrule
\end{tabular}
\end{table}

\subsection{Cross-Family Reference Model Validation}
\label{sec:cross_ref}

To test whether DDR is specific to the Pythia-410M reference model, we evaluate with two alternative references: OLMo-2-1B and Qwen3-0.6B (Table~\ref{tab:cross_ref}). Drift direction is consistent across all references (all drifting features show negative slopes), and shape features remain stationary. However, drift magnitude varies substantially: OLMo reference shows 51\% of the Pythia reference magnitude for entropy\_mean drift ($-0.61$ vs.\ $-1.19$). LOGO F1 is relatively stable across references (0.672--0.689), indicating that classification performance is robust even when drift magnitude varies.

\begin{table}[t]
\centering
\small
\caption{Cross-reference-model DDR comparison.}
\label{tab:cross_ref}
\begin{tabular}{lcccc}
\toprule
\textbf{Ref.\ Model} & \textbf{LOGO F1} & \textbf{ent\_mean $\rho$} & \textbf{hapax $\rho$} & \textbf{Drift mag.} \\
\midrule
Pythia-410M & 0.680 & $-0.60$ & $-0.63$ & 1.19 \\
OLMo-2-1B & 0.672 & $-0.42$ & $-0.58$ & 0.61 \\
Qwen3-0.6B & 0.689 & $-0.55$ & $-0.61$ & 0.95 \\
\bottomrule
\end{tabular}
\end{table}

\subsection{Downstream Utility (Negative Finding)}
\label{sec:downstream}

We fine-tune OLMo-1B on depth-stratified subsets to test whether depth-aware data curation improves downstream CLM quality (Table~\ref{tab:downstream}). Depth 0+1 filtered data (ppl $= 2.387$) performs comparably to size-matched random subsampling (ppl $= 2.367$). TOST confirms equivalence at $\pm 2\%$ margin ($p = 1.1 \times 10^{-4}$). The full dataset achieves the best perplexity (2.335), and depth-0-only data is substantially worse (2.681). Cross-seed variance is extremely low (std $< 0.002$), confirming the robustness of this negative result.

\begin{table}[t]
\centering
\small
\caption{Downstream CLM perplexity by training data depth composition (3 seeds).}
\label{tab:downstream}
\begin{tabular}{lcccc}
\toprule
\textbf{Variant} & \textbf{Size} & \textbf{PPL} & \textbf{$\Delta$ vs All} & \textbf{TOST $p$} \\
\midrule
All depths & 12,299 & 2.335$\pm$0.001 & --- & --- \\
Size-matched random & 3,951 & 2.367$\pm$0.001 & $+0.032$ & ref \\
Depth 0+1 only & 3,951 & 2.387$\pm$0.001 & $+0.052$ & $1.1\!\times\!10^{-4}$ \\
Binary filtered & 3,951 & 2.389$\pm$0.001 & $+0.054$ & $8.7\!\times\!10^{-5}$ \\
Depth 0 only & 797 & 2.681$\pm$0.001 & $+0.346$ & --- \\
\bottomrule
\end{tabular}
\end{table}

\subsection{Prompt-Type Comparison}
\label{sec:prompt}

Table~\ref{tab:prompt} compares three prompt types on Pythia-2.8B. Continuation chains yield the highest CV F1 (0.781), followed by inspired rewriting (0.737) and elaboration (0.730), with a cross-prompt-type range of 5.1 pp. Leave-one-prompt-type-out (LOPTO) evaluation shows consistent generalization (avg F1 = 0.752). Elaboration chains exhibit higher depth retention ($d_4/d_0$ ratio = 0.986 vs.\ 0.831 for continuation) but weaker entropy drift ($\rho = -0.471$ vs.\ $-0.557$).

\begin{table}[t]
\centering
\small
\caption{Per-prompt-type classification and drift statistics.}
\label{tab:prompt}
\begin{tabular}{lccccc}
\toprule
\textbf{Type} & \textbf{CV F1} & \textbf{LOPTO} & \textbf{Ret.} & \textbf{ent $\rho$} & \textbf{hpx $\rho$} \\
\midrule
Continuation & 0.781 & 0.762 & 0.831 & $-0.557$ & $-0.651$ \\
Elaboration & 0.730 & 0.733 & 0.986 & $-0.471$ & $-0.674$ \\
Rewriting & 0.737 & 0.762 & 0.939 & $-0.477$ & $-0.670$ \\
\bottomrule
\end{tabular}
\end{table}

\subsection{Paraphrastic vs.\ Non-Paraphrastic Chains}
\label{sec:para}

Table~\ref{tab:para} compares paraphrastic (semantic-preserving) and non-paraphrastic (continuation) chains. Paraphrastic chains show attractor convergence (depth retention = 0.995, effective classes $\approx 2$), confirming theoretical predictions. The depth signal is weaker in paraphrastic chains (binary F1 = 0.665 vs.\ 0.694, $\Delta = 2.9$ pp), and entropy drift is reduced ($\rho = -0.380$ vs.\ $-0.557$).

\begin{table}[t]
\centering
\small
\caption{Paraphrastic vs.\ non-paraphrastic chain comparison.}
\label{tab:para}
\begin{tabular}{lccccc}
\toprule
\textbf{Chain Type} & \textbf{CV F1} & \textbf{Bin.\ F1} & \textbf{Ret.} & \textbf{ent $\rho$} & \textbf{Eff.\ cls} \\
\midrule
Non-para.\ (cont.) & 0.781 & 0.694 & 0.831 & $-0.557$ & 3 \\
Paraphrastic & 0.723 & 0.665 & 0.995 & $-0.380$ & $\sim$2 \\
\bottomrule
\end{tabular}
\end{table}

\subsection{Cross-Model Chains}
\label{sec:cross_model}

Cross-model chains (different generator at each depth step) simulate a more realistic web contamination scenario where multiple models contribute to a generation chain. As shown in Table~\ref{tab:cross_model}, cross-model chains yield substantially higher binary F1 (0.764 vs.\ 0.694, $+7.0$ pp) and LOGO binary F1 (0.830 vs.\ 0.694), because generator heterogeneity amplifies the distributional shift at each depth step.

\begin{table}[t]
\centering
\small
\caption{Cross-model vs.\ single-model chain classification.}
\label{tab:cross_model}
\begin{tabular}{lccc}
\toprule
\textbf{Setting} & \textbf{Binary F1} & \textbf{3-class F1} & \textbf{LOGO Bin.} \\
\midrule
Single-model & 0.694 & 0.762 & 0.694 \\
Cross-model & 0.764 & 0.789 & 0.830 \\
$\Delta$ & $+7.0$ pp & $+2.7$ pp & $+13.6$ pp \\
\bottomrule
\end{tabular}
\end{table}

\subsection{Multi-Domain Seeds}
\label{sec:domain}

Table~\ref{tab:domain} shows per-domain binary F1 across four seed domains. Wikipedia yields the highest F1 (0.724), while scientific text yields the lowest (0.631), a gap of 9.3 pp. Cross-domain LOGO (train on 3 domains, test on 1) yields 0.652.

\begin{table}[t]
\centering
\small
\caption{Per-domain binary F1 (LOGO across generators, within-domain).}
\label{tab:domain}
\begin{tabular}{lcccc}
\toprule
\textbf{Domain} & \textbf{Seeds} & \textbf{Binary F1} & \textbf{95\% CI} & \textbf{$\Delta$ vs Wiki} \\
\midrule
Wikipedia & 500 & 0.724 & [0.704, 0.744] & --- \\
News & 500 & 0.667 & [0.645, 0.689] & $-5.7$ pp \\
Creative & 500 & 0.658 & [0.635, 0.681] & $-6.6$ pp \\
Scientific & 500 & 0.631 & [0.607, 0.655] & $-9.3$ pp \\
\midrule
\textbf{Cross-domain} & 2000 & \textbf{0.652} & [0.638, 0.666] & --- \\
\bottomrule
\end{tabular}
\end{table}

\subsection{Finer-Grained Ordinal Classification}
\label{sec:ordinal}

We test whether the depth signal supports finer resolution beyond 3 classes (Table~\ref{tab:ordinal}). Five-class macro F1 drops to 0.468, with ordinal methods (CORAL: 0.472, Frank-Hall: 0.465) offering no meaningful improvement. Confusion rates between adjacent deep classes are high: gen-1/gen-2 (43\%), gen-2/gen-3 (51\%), gen-3/gen-4 (48\%). Mean absolute depth error is 0.89. This confirms that 3-class resolution represents a practical ceiling imposed by signal strength.

\begin{table}[t]
\centering
\small
\caption{Ordinal classification resolution comparison.}
\label{tab:ordinal}
\begin{tabular}{llcc}
\toprule
\textbf{Granularity} & \textbf{Method} & \textbf{Macro F1} & \textbf{Ordinal $\Delta$} \\
\midrule
3-class & XGBoost & 0.762 & --- \\
5-class & XGBoost & 0.468 & --- \\
5-class & CORAL & 0.472 & $+0.4$ pp \\
5-class & Frank-Hall & 0.465 & $-0.3$ pp \\
\bottomrule
\end{tabular}
\end{table}

\subsection{Human Edit Robustness}
\label{sec:robustness}

Table~\ref{tab:robust} reports robustness to simulated human edits. Even 5\% random word replacement reduces binary F1 by 4.6 pp (0.694 $\to$ 0.648), and LLM-simulated rephrasing (edit distance 0.51) reduces it by 5.3 pp (0.694 $\to$ 0.641). At 20\% replacement, binary F1 drops to 0.583.

\begin{table}[t]
\centering
\small
\caption{Robustness to simulated human edits.}
\label{tab:robust}
\begin{tabular}{lccc}
\toprule
\textbf{Perturbation} & \textbf{Binary F1} & \textbf{$\Delta$} & \textbf{Edit Dist.} \\
\midrule
None (baseline) & 0.694 & --- & 0.0 \\
Random 5\% & 0.648 & $-4.6$ pp & $\sim$0.05 \\
Random 10\% & 0.611 & $-8.3$ pp & $\sim$0.10 \\
Random 20\% & 0.583 & $-11.1$ pp & $\sim$0.20 \\
LLM rephrase & 0.641 & $-5.3$ pp & 0.51 \\
\bottomrule
\end{tabular}
\end{table}

\subsection{LD-Score as Additional Feature}

Adding LD-Score (character-distribution signature) as an 8th feature provides marginal improvement: macro F1 increases from 0.773 to 0.779 ($+0.6$ pp). LD-Score alone achieves 0.623 F1, with Spearman $\rho = -0.31$ with depth and feature importance rank of 0.028 (lowest). The character-level signal is largely subsumed by existing word-level features.

\subsection{End-to-End Neural Classifier}
\label{sec:e2e}

Table~\ref{tab:e2e} compares the 7-feature ensemble against a QLoRA fine-tuned OLMo-1B classifier. Within-distribution, QLoRA substantially outperforms (binary F1: 0.842 vs.\ 0.694, $+14.8$ pp; 3-class: 0.891 vs.\ 0.762, $+12.9$ pp). Under LOGO, the gap narrows: QLoRA 3-class F1 drops to 0.743 (below 7-feature's 0.762), while QLoRA binary F1 remains higher (0.712 vs.\ 0.694). QLoRA incurs 47$\times$ inference cost. For 5-class, QLoRA achieves 0.571 vs.\ 0.468 for XGBoost.

\begin{table}[t]
\centering
\small
\caption{End-to-end (QLoRA) vs.\ 7-feature ensemble.}
\label{tab:e2e}
\begin{tabular}{lccccc}
\toprule
\textbf{Method} & \textbf{3-cl} & \textbf{Bin} & \textbf{LOGO 3} & \textbf{LOGO Bin} & \textbf{Cost} \\
\midrule
7-feat XGB & .762 & .694 & .762 & .694 & 1$\times$ \\
QLoRA & .891 & .842 & .743 & .712 & 47$\times$ \\
\bottomrule
\end{tabular}
\end{table}

\subsection{Text Length Control}
\label{sec:length}

Text length correlates weakly with depth (Spearman $\rho = 0.12$, $p = 0.001$). Length-matched subsampling reduces binary F1 from 0.770 to 0.724 ($-4.6$ pp), confirming that length contributes to but does not primarily drive the depth signal. Length-normalized features preserve depth trends for both TTR and hapax ratio.

\begin{table}[t]
\centering
\small
\caption{Length control analysis.}
\label{tab:length}
\begin{tabular}{lcc}
\toprule
\textbf{Condition} & \textbf{Binary F1} & \textbf{$\Delta$} \\
\midrule
Unmatched (original) & 0.770 & --- \\
Length-matched subsample & 0.724 & $-4.6$ pp \\
Length-normalized features & 0.748 & $-2.2$ pp \\
\bottomrule
\end{tabular}
\end{table}

\subsection{Scale Effect: Larger Generator}
\label{sec:scale}

Adding Qwen3-8B-Base (Table~\ref{tab:scale}) reveals that depth signal degrades at larger scale. Qwen3-8B achieves LOGO binary F1 = 0.637, 10.6 pp below the 6-generator average. Feature drift magnitudes compress: hapax $\rho$ drops from average $-0.63$ to $-0.41$, and entropy\_mean $\rho$ from $-0.60$ to $-0.33$ (compression ratio 0.58). Within-CV F1 remains reasonable (0.711), but cross-generator transfer to 8B is substantially weaker.

\begin{table}[t]
\centering
\small
\caption{Scale effect: Qwen3-8B vs.\ smaller generators.}
\label{tab:scale}
\begin{tabular}{lcccc}
\toprule
\textbf{Generator} & \textbf{Size} & \textbf{LOGO Bin.} & \textbf{hapax $\rho$} & \textbf{ent $\rho$} \\
\midrule
Pythia-1B & 1.0B & 0.724 & $-0.65$ & $-0.60$ \\
OLMo-2-1B & 1.5B & 0.643 & $-0.58$ & $-0.55$ \\
Pythia-2.8B & 2.8B & 0.711 & $-0.63$ & $-0.60$ \\
Qwen3-8B & 8.0B & 0.637 & $-0.41$ & $-0.33$ \\
\bottomrule
\end{tabular}
\end{table}

\subsection{Ordinal Baseline Comparison}
\label{sec:baselines}

Table~\ref{tab:baselines} compares the 7-feature ensemble against existing detector scores repurposed as ordinal depth predictors. The 7-feature ensemble (LOGO binary F1 = 0.680) outperforms all single-detector baselines, but the margin is modest: $+2.7$ pp vs.\ Avg-LogProb (bootstrap $p = 0.038$), $+3.9$ pp vs.\ LD-Score, $+6.8$ pp vs.\ FastDetectGPT, $+10.9$ pp vs.\ Binoculars. A multi-detector ensemble (4 detectors + surface features) achieves 0.708, surpassing the 7-feature ensemble by 2.8 pp ($p = 0.024$).

\begin{table}[t]
\centering
\small
\caption{Ordinal depth classification: 7-feature vs.\ existing detector baselines.}
\label{tab:baselines}
\begin{tabular}{lcccc}
\toprule
\textbf{Method} & \textbf{LOGO Bin F1} & \textbf{95\% CI} & \textbf{$\Delta$} & \textbf{$p$} \\
\midrule
\textbf{7-feature ens.} & \textbf{0.680} & [.672, .688] & --- & --- \\
Avg-LogProb & 0.653 & [.612, .694] & $-2.7$ & .038 \\
LD-Score & 0.641 & [.598, .684] & $-3.9$ & .011 \\
FastDetectGPT & 0.612 & [.571, .653] & $-6.8$ & .002 \\
Binoculars & 0.571 & [.528, .614] & $-10.9$ & $<.001$ \\
\midrule
4-det + surface & 0.708 & [.694, .722] & $+2.8$ & .024 \\
\bottomrule
\end{tabular}
\end{table}

%=====================================================================
\section{Discussion}

Our experiments reveal several important patterns and limitations:

\paragraph{DDR is real but modest.} The differential depth response---drifting vs.\ stationary feature behavior---is consistently observed across generators, reference models, and prompt types. However, the practical signal strength is limited: LOGO binary F1 = 0.680 represents moderate discriminability, not reliable classification.

\paragraph{Surface features dominate.} Surface-only features achieve 95.9\% of the full feature set's performance (0.741 vs.\ 0.773), suggesting that reference features provide complementary but marginal signal. This raises the question of whether the reference model computation (requiring a 410M model forward pass) is justified for 3.2 pp improvement.

\paragraph{Scope limitations compound.} The depth signal is restricted to base models, moderate temperatures, and unedited text. Each condition further reduces performance: instruct models collapse the signal; larger models compress feature margins; human editing degrades F1 by 5+ pp. These limitations compound to define an increasingly narrow practical scope.

\paragraph{Negative downstream finding is robust.} Depth-aware curation is statistically equivalent to random subsampling across all tested conditions. This establishes depth fingerprinting as a diagnostic tool---useful for provenance characterization and contamination auditing---but not a curative tool for data selection.

\paragraph{Existing detectors provide comparable signal.} The multi-detector ensemble surpasses the 7-feature ensemble by 2.8 pp, suggesting that repurposing existing detection scores may be a more practical approach to depth estimation than designing depth-specific features.

%=====================================================================
\section{Conclusion}

We formulated generation depth classification as a task beyond binary AI-text detection and characterized differential depth response (DDR) across 5 base-model generators from 3 families. A 7-feature distributional ensemble achieves LOGO binary F1 = 0.680 with a single 410M reference model. Cross-model chains yield binary F1 = 0.764. Depth-aware curation does not improve downstream quality (TOST confirmed). The depth signal degrades with model scale (8B: $-10.6$ pp), human editing ($-5.3$ pp for light rephrasing), and domain shift (scientific text: $-9.3$ pp vs.\ Wikipedia). These scope boundaries---base models, moderate temperatures, unedited text---define the current applicability of depth fingerprinting.

Future work should address the instruct-model gap, scale to larger generators, and investigate whether DDR-derived features can complement existing detection methods as auxiliary signals rather than standalone depth classifiers.

%=====================================================================
\section*{Limitations}

This study has several important limitations. (1) All experiments use base (non-instruction-tuned) models; instruction tuning substantially degrades the depth signal. (2) The generator scale tested (0.6B--8B) excludes frontier models (70B+), which likely compress feature margins further. (3) Human edit robustness is tested with simulated edits only. (4) Temperature sensitivity restricts reliable performance to $t \in [0.5, 0.9]$. (5) The 3-class formulation represents a ceiling: finer-grained depth estimation (5-class) achieves only 0.468 F1. (6) The multi-detector ensemble outperforms the proposed features, questioning the marginal contribution of depth-specific feature engineering. (7) All generation chains are constructed in controlled settings; in-the-wild chains may differ substantially.

\bibliography{custom}

\end{document}
